\section{Introduction}
\label{sec:introduction}
A failed attempt tells a language agent that something went wrong. It does
not tell the agent what to learn. In a long task, failure may follow an
unproductive search, a missed change of subgoal, or an earlier error from
which the agent never recovered. A research model can inspect the
trajectory and suggest a correction. Reflection-based self-training uses
such feedback to improve generated samples \citep{dou2024rerest}.
Yet a convincing suggestion is still a hypothesis
about the policy's future behavior. Opening a cabinet instead of continuing
to search might reveal the needed object, or it might spend the remaining
budget and leave the task unsolved. The useful question is therefore concrete:
what happens when \emph{this policy} takes the proposed change from
\emph{that state}?

This question sits at a critical handoff in agent self-improvement. Memory
and harness adaptation can preserve experience outside the model
\citep{shinn2023reflexion,yu2026recuris}; policy training can carry useful
experience into the model parameters themselves
\citep{chen2026evotrainer,yang2026frontis}. The object we ultimately want to
retain is a policy that performs better when these external aids are
removed. A stronger intervention can improve the assisted trajectory
while leaving that object unchanged. Aspire shows why completing the
update loop is insufficient: agents can execute training while hidden-evaluation
gains remain sparse or unstable \citep{wu2026aspire}. We study the smaller,
more testable step between these two ends---how evidence from a failed attempt
becomes supervision that still helps after the research process is gone.

Prior work already supplies important pieces. CSO tests alternative actions
through policy continuation and learns from verified preferences
\citep{li2026cso}. SIRI validates skill-guided rollouts and distills useful
actions into an unassisted policy \citep{he2026siri}. OPID and SEED turn
trajectory hindsight into skills whose effect on action probabilities guides
policy learning \citep{yang2026opid,wu2026seed}. These advances make two
questions especially salient. First, which failure-derived changes actually
deserve to enter training? Second, once a change has earned that evidence,
what information should the learner receive from it?

The first question requires a baseline for the counterfactual. Restarting a
failed attempt gives the policy another opportunity to succeed even without
the proposed change. RePair therefore restores the source decision state and
runs two matched continuations: the registered baseline action and the
candidate intervention, with the same frozen policy available for subsequent
decisions. A bounded intervention can itself complete the task; its success
then measures the intervention's utility, not autonomous policy behavior.
The terminal outcomes measure whether adding the intervention changes
completion from that point. This same-state comparison
turns a plausible correction into experimentally grounded evidence.

The second question begins after verification. An action target teaches the
learner what to do at the source state. A compact strategy target can also
teach the active subgoal, the progress signal to watch for, and the condition
for changing course. In the cabinet example, opening the door is one action;
recognizing what observation would complete the search, and how to recover if
the object is absent, describes the decision structure around that action.
The two targets define a controlled comparison on the same verified states
and corrected actions. Here we instantiate the dual-view update and evaluate
its descendant through the original action-only interface, with external
failure memory and repair intervention disabled. An action-only training
control remains necessary to isolate the auxiliary view's contribution.

We introduce \method{} to make this handoff measurable end to end
(Figure~\ref{fig:framework}). Failure analysis and cross-round experience
help propose testable changes; same-state continuations determine which ones
become positive evidence; action and strategy targets specify
what enters learning; scaffold-free evaluation measures what remains in the
policy. The surrounding Analyzer, Planner, and
memory machinery supports this sequence, while the empirical study tracks
which of these handoffs actually succeeds.

Our ALFWorld study starts from Qwen2.5-3B-Instruct. Across five development
rounds, most verified interventions leave both branches unsuccessful.
Trace analysis locates a recurring failure: a local correction executes,
but the continuation policy misses the next prerequisite or repeats an
already completed action. A later validation makes more of the proposed
strategy executable and obtains stable terminal benefits. Yet training on
four such sources does not improve full-task selection success. These
observations locate the unresolved handoff: a program can finish a task
while the training examples teach only its initial decision. We report
the trace and supervision evidence behind this gap and measure its
consequence for the policy after external assistance is removed.
